% Appendix: EA400 Binary File Format Specification
% Reference for the reverse-engineered .BIN format.

\appendix
\section{EA400 echosounder binary file format}
\label{app:bin_format}

The Echologger EA400 echosounder stores data in a proprietary binary
format (.BIN) for which no public documentation or parsing tools exist
from the manufacturer. The file format was reverse-engineered from hex
analysis of instrument files across three deployment sites, verified
against concurrent ASCII text (.log) outputs for the same deployments.

\subsection{File structure}

Each .BIN file consists of a 128-byte file header followed by a
sequence of fixed-size records, one per acoustic ping. The record size
depends on the configured number of depth bins $N$:
$R = 64 + 2N + 64$~bytes.

\subsection{File header (128 bytes)}

\begin{tabular}{llll}
\hline
Offset & Type & Field & Notes \\
\hline
0x00 & uint32 LE & header\_size & Always 128 \\
0x08 & char[24] & timer\_name & ``Default timer'' \\
0x3C & float32 LE & sample\_interval\_s & Typically 0.5 \\
0x58 & float32 LE & sound\_speed\_mps & Typically 1500 \\
\hline
\end{tabular}

\subsection{Data record}

Each record contains two sub-records:

\subsubsection{DATA sub-record (64 + 2N bytes)}

\begin{tabular}{llll}
\hline
Offset & Type & Field & Notes \\
\hline
+0 & char[4] & marker & ``DATA'' (0x44415441) \\
+4 & uint16 LE & version & 0x0101 \\
+6 & uint16 LE & padding & \\
+8 & uint32 LE & ping\_number & 1-indexed \\
+12 & uint16 LE & num\_samples & $N$ (272 or 400) \\
+14 & uint16 LE & padding & \\
+16 & uint32 LE & timestamp & Unix epoch seconds$^*$ \\
+20 & uint32 LE & subsecond & Internal counter \\
+24--63 & mixed & prev\_metadata & Echo of previous STAT; ignore \\
+64 & uint16[$N$] LE & backscatter & 0--1023 per depth bin \\
\hline
\end{tabular}

$^*$The timestamp represents instrument local time encoded as Unix
epoch seconds (i.e., the local date/time is treated as if it were UTC
when computing the epoch value). A timezone offset must be applied
during post-processing.

\subsubsection{STAT sub-record (64 bytes)}

The STAT sub-record contains the authoritative per-ping metadata.
Values in the DATA sub-record header (positions +24--63) are echoes
of the \emph{previous} ping's STAT and should not be used.

\begin{tabular}{llll}
\hline
Offset & Type & Field & Units \\
\hline
+0 & uint16 LE & marker & 0x0102 \\
+2 & uint16 LE & padding & \\
+4 & uint32 LE & ping\_number & \\
+8 & uint32 LE & num\_samples & \\
+12 & float32 LE & temperature & \si{\celsius} \\
+16 & float32 LE & battery\_V & \si{\volt} (tentative) \\
+20 & float32 LE & altitude & \si{\meter} (0 = no echo) \\
+24 & float32 LE & pitch & \si{\degree} \\
+28 & float32 LE & roll & \si{\degree} \\
+32--63 & zero & reserved & \\
\hline
\end{tabular}

\subsection{Observed configurations}

\begin{tabular}{lccc}
\hline
Site & $N$ (depth bins) & Record size & Bin spacing \\
\hline
SouthSIOPier (MOP511) & 272 & 672 bytes & \SI{7.5}{\milli\meter} \\
TorreyPines (MOP586) & 400 & 928 bytes & \SI{7.5}{\milli\meter} \\
SolanaBeach (MOP654) & 272 & 672 bytes & \SI{7.5}{\milli\meter} \\
\hline
\end{tabular}

Typical file sizes range from \SI{70}{\mega\byte} (short test
deployments, $\sim$100 pings) to \SI{525}{\mega\byte} (full
deployments, $\sim$780{,}000 pings at \SI{2}{\hertz} burst sampling).
All multi-byte values are little-endian.
